\section{7. What the Five Functions Hide: Specification Difficulty and Verification Strength}

The five-function lens describes \emph{what} a system does without surfacing \emph{what makes each function hard}. Two variables, hidden inside that view, drive how much autonomy a system can safely be given.

\textbf{The first is specification difficulty}, inside context assembly. This sits on a gradient: \textbf{Formal} (fully specifiable in structured form --- an API schema, a complete checklist); \textbf{Expert-articulable} (a domain expert can explain the requirement, but doing so takes real professional knowledge, and a non-expert can't assess whether the specification is complete); \textbf{Tacit} (the expert "knows it when they see it" but can't fully pre-specify it --- strategic judgment, reading a counterparty).

The tacit end of this gradient is grounded in Polanyi's Paradox --- "we know more than we can tell" --- which is an epistemic claim about people: that they can know things they cannot fully articulate. It does not, by itself, establish that a machine-learning system could never approximate that knowledge through some other route --- demonstration, outcome-based learning, extended interaction. \textbf{The irreducibility of the human contribution at the tacit end of this gradient is this framework's own governance-relevant working assumption, not a conclusion Polanyi's work hands us directly.} Whether tacit expertise genuinely resists capture by any learning method, or merely resists capture by current methods and current data, is a live and unresolved question, returned to in Section 15.

\textbf{The second hidden variable is verification strength}, inside the verification function. "Oracle" is used here as a general umbrella term for "the standard against which an output can be checked" --- but the term can mean at least five materially different things, and it's worth distinguishing them explicitly rather than letting them collapse together: a deterministic automated test, a measured empirical outcome, an expert reviewer's judgment, an institution empowered to authoritatively resolve a question (like a court), and unresolved ground truth itself. A court is an \textbf{authoritative decision-maker} for a legal question without necessarily being the only source of evidence about whether a particular argument is doctrinally sound; an expert reviewer is a \textbf{fallible evaluator}, not an Oracle in the same sense a compiler is. The gradient underneath the umbrella term has five rungs:

\begin{itemize}
\item \textbf{Formal validator} --- a deterministic rule or executable test (a compiler, a schema check). Verification cost is near zero, reliability is high.
\item \textbf{Empirical evaluator} --- a statistical or measurable check with residual uncertainty (a benchmark score, a tolerance measurement).
\item \textbf{Expert assessor} --- evaluation requires human professional judgment; throughput-limited, imperfect inter-rater agreement even among qualified experts.
\item \textbf{Authoritative decision-maker} --- an institution empowered to resolve the question even where the underlying reasoning is genuinely contested (a court's ruling is authoritative regardless of whether every lawyer agrees with the reasoning).
\item \textbf{Unresolved ground truth} --- the outcome simply isn't known yet, and no institution can resolve it in advance (a negotiation's actual outcome, a market's future move).
\end{itemize}

This sharper breakdown separates authority from correctness --- a distinction worth making explicit, since it resolves a common conflation in discussions of "scaling optimism" (the view that a large enough model will eventually internalize its own verification standard): no volume of training data changes where the \emph{institutional authority} to determine legal correctness resides, but that, on its own, doesn't settle whether a model can improve at \textbf{predicting} a court's likely decision, \textbf{identifying} controlling authority, \textbf{recognizing} standard doctrinal errors, or \textbf{producing analysis domain experts rate as sound}. These are meaningfully different capabilities --- normative authority, doctrinal accuracy, and predictive accuracy are not the same claim, and a model can plausibly improve at the latter two without that touching the first at all.

\textbf{The resulting task space, plotting specification difficulty against verification strength, produces four characteristic regions:}

\emph{Low specification difficulty, strong verification --- "automate and verify."} Both the specification and the check are available up front; the system can operate with real autonomy because errors are both inexpensive to produce and cheap to detect.

\emph{Low specification difficulty, weak verification --- "fluent output, no feedback signal."} The least well-understood region and the one most prone to misallocated confidence: the inputs are simple, the outputs are fluent, and the system appears to run without friction, but fluency substitutes for quality assessment when there's nothing to verify against, and errors accumulate silently.

\emph{High specification difficulty, strong verification --- "augment and verify."} An expert does the hard work of framing the task, and a checklist, playbook, or professional standard provides a genuine verification mechanism at the end. The highest-value target for most professional-services deployment.

\emph{High specification difficulty, weak verification --- "expert judgment throughout."} The human contribution dominates every function. The system can contribute to sub-tasks --- enumerating scenarios, surfacing options --- but can't close the loop end to end.

Task structure genuinely constrains the \emph{maximum defensible} autonomy for a task, but \textbf{actual} autonomy in a given deployment also depends on system architecture, error tolerance, whether the relevant actions are reversible or irreversible, an organization's risk appetite, the monitoring and incident-response capability actually in place, and legal or professional obligations specific to the context. An organization can always choose to deploy \emph{less} autonomy than a task might technically support, and often should. The task space tells you the ceiling, not the number an organization is obligated to hit.

\textbf{Supporting and adjacent literature.} Kahneman and Klein's 2009 synthesis of the heuristics-and-biases and naturalistic-decision-making traditions identified two conditions for reliable expert intuition: a high-validity environment with learnable, stable cue-outcome relationships, and adequate opportunity to learn those regularities through fast, clear feedback. By analogy --- not as their stated finding --- a high-validity environment resembles a strong Oracle in this framework's sense: the domain provides clear, timely feedback on whether a judgment was right. It's worth noting, since it complicates a simple mapping, that Kahneman and Klein themselves also observed that \textbf{consistent algorithms can outperform human judgment specifically in low-validity or noisy environments}, precisely because algorithms don't suffer the inconsistency that noisy human judgment does under weak cues --- a point that cuts against any simple claim that "if human judgment fails here, no automated method can succeed here either." Their framework is a productive analogy for thinking about verification strength, not a theory of automated verification in its own right.

Christiano's "easy to verify, hard to generate" asymmetry from the AI-alignment literature formalizes a related intuition: delegation works when verification cost is low relative to generation cost. Coding agents succeed partly because compilers and test suites provide a cheap, strong check, not necessarily because models are better at code than at prose. This report uses the term \textbf{"verification tax"} for the specific idea that fine-tuning can raise a model's generation confidence without strengthening the underlying verification standard --- that phrase is this framework's own term for the concept, not a term of art with a single settled meaning elsewhere in the literature, where similar phrasing is sometimes used for related but distinct costs, such as the statistical cost of auditing model calibration.

The Cynefin framework's classification of decision contexts as clear, complicated, complex, or chaotic is organized around the legibility of cause-and-effect relationships and applicable constraints in a domain. The mapping from Cynefin onto verification strength offered here is this report's own comparison, and an imperfect one --- Cynefin was built to classify decision contexts generally, not specifically to model where an automated verification standard exists.

Sutton's "Bitter Lesson" --- that general methods leveraging computation, especially search and learning, have historically outperformed systems built on hand-encoded domain knowledge --- is broader than a claim about within-model scaling specifically; it's a claim about method design across the history of AI (chess, Go, speech recognition, computer vision). Even taken at full strength, it doesn't by itself resolve questions about external validity, institutional authority, or verification --- it's a claim about which methods improve with more compute and data, not a claim that improving methods this way closes the gap between a model's output and independently-checkable ground truth.

Parasuraman, Sheridan, and Wickens's four-stage model of human-automation interaction (information acquisition, information analysis, decision/action selection, action implementation, each with its own appropriate level of automation) is a real, independently established framework. The connection to "Oracle strength" drawn here is this report's extension of their model --- their own argument is that appropriate automation level depends on failure consequences and automation reliability, which is adjacent to but not identical with the verification-strength argument made in this section.

\textbf{Three counter-arguments, addressed directly.} \emph{Scaling optimism} --- that a sufficiently capable model will eventually internalize the verification standard itself --- runs into the authority/accuracy distinction drawn above: it may be right about predictive or doctrinal accuracy improving with scale, while still being wrong to imply that institutional authority migrates to the model as a result. These are separable claims, and conflating them is the actual error in the optimistic view, not scaling itself. \emph{The Bitter Lesson} doesn't resolve this either way --- it's a claim about method design, not about where verification authority resides. \emph{Agentic-loop compensation} --- that tool use and reflection can compensate for a weak Oracle --- has genuine merit where the tools invoked have strong verification of their own (a calculator, a compiler, a cryptographic signature check, a live database with a defined schema): these can genuinely terminate a verification regress rather than merely relocating it, because they supply independently validated evidence rather than another layer of unverified judgment. The regress problem is real specifically where each tool's output itself depends on another equally ungrounded judgment --- it is not an unconditional wall that agentic tool use always hits, and the distinction between tools that terminate the regress and tools that merely relocate it is worth checking case by case.
